\documentclass[12pt,a4paper]{article}

% ---- Encoding & Font ----
\usepackage{iftex}
\ifPDFTeX
  \usepackage[utf8]{inputenc}
  \usepackage[T1]{fontenc}
\else
  \usepackage{fontspec}  % XeTeX/LuaTeX (tectonic): Unicode nativo (·, —, ¿, ª, §)
\fi
\usepackage[spanish]{babel}
\usepackage{csquotes}

% ---- Page layout ----
\usepackage[top=2.5cm, bottom=2.5cm, left=2.5cm, right=2.5cm]{geometry}
\usepackage{setspace}
\onehalfspacing

% ---- Math ----
\usepackage{amsmath, amssymb, amsthm}
\usepackage{mathtools}
\usepackage{physics}
\usepackage{esint}  % \oiint

% ---- Graphics ----
\usepackage{graphicx}
\usepackage{tikz}
\usepackage{float}
\usepackage{caption}

% ---- Tables ----
\usepackage{array}
\usepackage{booktabs}
\usepackage{tabularx}
\usepackage{colortbl}

% ---- Colours ----
\usepackage{xcolor}
\definecolor{notebg}{HTML}{FFF3CD}
\definecolor{noteborder}{HTML}{FFC107}
\definecolor{boxred}{HTML}{F8D7DA}
\definecolor{boxredborder}{HTML}{F5C6CB}

% ---- Boxes ----
\usepackage{mdframed}
\usepackage{tcolorbox}
\tcbuselibrary{most}

\newtcolorbox{keybox}[1][]{
  colback=blue!5,
  colframe=blue!70!black,
  arc=4pt,
  boxrule=1pt,
  fonttitle=\bfseries,
  title={#1}
}

\newtcolorbox{warnbox}[1][]{
  colback=boxred,
  colframe=boxredborder,
  coltitle=black!75,
  arc=4pt,
  boxrule=1pt,
  fonttitle=\bfseries,
  title={#1}
}

\newtcolorbox{notebox}[1][]{
  colback=notebg,
  colframe=noteborder,
  coltitle=black!75,
  arc=4pt,
  boxrule=1pt,
  fonttitle=\bfseries,
  title={#1}
}

% ---- Diagramas (gráficas y esquemas) ----
\usepackage{pgfplots}
\pgfplotsset{compat=1.18}
\usetikzlibrary{babel}  % babel español activa `>`; esto evita romper las flechas `->`
% courses/ElecMag2024/Clases/assets/tikzstyles.tex
% ---------------------------------------------------------------------------
% Estilos compartidos para los diagramas del curso Electromagnetismo (2024).
% Fuente única del "look" visual (colores, grosores, escalas) para que todas
% las figuras (TikZ / pgfplots / CircuiTikz) sean consistentes entre clases.
%
% Es una COPIA de courses/Fisica3-2015/Clases/assets/tikzstyles.tex: mismo
% dominio, misma paleta. Copia y no symlink para que cada curso sea
% autocontenido y la paleta de Física III pueda evolucionar aparte.
%
% Uso: la clase debe cargar antes `tikz` y `pgfplots` (y `circuitikz` si la
%      clase tiene circuitos), y luego % courses/ElecMag2024/Clases/assets/tikzstyles.tex
% ---------------------------------------------------------------------------
% Estilos compartidos para los diagramas del curso Electromagnetismo (2024).
% Fuente única del "look" visual (colores, grosores, escalas) para que todas
% las figuras (TikZ / pgfplots / CircuiTikz) sean consistentes entre clases.
%
% Es una COPIA de courses/Fisica3-2015/Clases/assets/tikzstyles.tex: mismo
% dominio, misma paleta. Copia y no symlink para que cada curso sea
% autocontenido y la paleta de Física III pueda evolucionar aparte.
%
% Uso: la clase debe cargar antes `tikz` y `pgfplots` (y `circuitikz` si la
%      clase tiene circuitos), y luego % courses/ElecMag2024/Clases/assets/tikzstyles.tex
% ---------------------------------------------------------------------------
% Estilos compartidos para los diagramas del curso Electromagnetismo (2024).
% Fuente única del "look" visual (colores, grosores, escalas) para que todas
% las figuras (TikZ / pgfplots / CircuiTikz) sean consistentes entre clases.
%
% Es una COPIA de courses/Fisica3-2015/Clases/assets/tikzstyles.tex: mismo
% dominio, misma paleta. Copia y no symlink para que cada curso sea
% autocontenido y la paleta de Física III pueda evolucionar aparte.
%
% Uso: la clase debe cargar antes `tikz` y `pgfplots` (y `circuitikz` si la
%      clase tiene circuitos), y luego \input{../assets/tikzstyles.tex}
% (compilar desde el directorio de la clase para que la ruta relativa resuelva).
% ---------------------------------------------------------------------------

% ---- Paleta (alineada con las cajas del preámbulo: keybox/notebox/warnbox) ----
\definecolor{figblue}{HTML}{1F4E79}   % azul principal (líneas, keybox)
\definecolor{figred}{HTML}{C0392B}    % rojo de énfasis (warnbox)
\definecolor{figamber}{HTML}{B8860B}  % ámbar (notebox)
\definecolor{figgray}{HTML}{555555}   % auxiliares, ejes, guías

% ---- CircuiTikz: escala y grosor de los componentes ----
% Guarda \ifdefined: la electrostática (Clases 1-14) no carga `circuitikz`, y
% sin la guarda el \input revienta con `Undefined control sequence \ctikzset`.
% Cuando el curso llegue a circuitos, el bloque se activa solo.
\ifdefined\ctikzset
  \ctikzset{
    bipoles/thickness=1,
    resistors/scale=0.9,
    capacitors/scale=0.9,
  }
\fi

% ---- TikZ: estilos reutilizables ----
% Nota: las puntas se declaran como `-{Latex[...]}` y no como `->`. La forma
% con llaves NO contiene un `>` literal, así que es inmune al `>` activo que
% introduce babel español (ver CLAUDE.md §7.3.3).
\usetikzlibrary{arrows.meta,calc,angles,quotes}

\tikzset{
  figline/.style={figblue, line width=0.9pt},
  figaux/.style={figgray, dashed, line width=0.6pt},
  % vectores
  figvec/.style={figblue, line width=0.9pt, -{Latex[length=2.2mm,width=1.6mm]}},
  figvecres/.style={figred, line width=1.3pt, -{Latex[length=2.6mm,width=1.9mm]}},
  figvecaux/.style={figgray, line width=0.7pt, -{Latex[length=1.8mm,width=1.3mm]}},
  figaxis/.style={figgray, line width=0.6pt, -{Latex[length=2mm,width=1.4mm]}},
  % cargas puntuales
  figcarga/.style={circle, inner sep=0pt, minimum size=5.4pt, draw=none},
  figpos/.style={figcarga, fill=figred},
  figneg/.style={figcarga, fill=figblue},
  figneutra/.style={figcarga, fill=figgray},
  % etiquetas
  figlbl/.style={font=\small},
  figlblaux/.style={font=\footnotesize, figgray},
}

% ---- pgfplots: estilo base de las gráficas ----
\pgfplotsset{
  emfig/.style={
    width=11cm, height=6.5cm,
    axis lines=left,
    xlabel style={font=\small},
    ylabel style={font=\small},
    tick label style={font=\footnotesize},
    every axis plot/.append style={line width=1pt},
    grid=major,
    grid style={figgray!20},
    legend cell align=left,
    legend style={font=\footnotesize, draw=figgray!40},
  },
}

% (compilar desde el directorio de la clase para que la ruta relativa resuelva).
% ---------------------------------------------------------------------------

% ---- Paleta (alineada con las cajas del preámbulo: keybox/notebox/warnbox) ----
\definecolor{figblue}{HTML}{1F4E79}   % azul principal (líneas, keybox)
\definecolor{figred}{HTML}{C0392B}    % rojo de énfasis (warnbox)
\definecolor{figamber}{HTML}{B8860B}  % ámbar (notebox)
\definecolor{figgray}{HTML}{555555}   % auxiliares, ejes, guías

% ---- CircuiTikz: escala y grosor de los componentes ----
% Guarda \ifdefined: la electrostática (Clases 1-14) no carga `circuitikz`, y
% sin la guarda el \input revienta con `Undefined control sequence \ctikzset`.
% Cuando el curso llegue a circuitos, el bloque se activa solo.
\ifdefined\ctikzset
  \ctikzset{
    bipoles/thickness=1,
    resistors/scale=0.9,
    capacitors/scale=0.9,
  }
\fi

% ---- TikZ: estilos reutilizables ----
% Nota: las puntas se declaran como `-{Latex[...]}` y no como `->`. La forma
% con llaves NO contiene un `>` literal, así que es inmune al `>` activo que
% introduce babel español (ver CLAUDE.md §7.3.3).
\usetikzlibrary{arrows.meta,calc,angles,quotes}

\tikzset{
  figline/.style={figblue, line width=0.9pt},
  figaux/.style={figgray, dashed, line width=0.6pt},
  % vectores
  figvec/.style={figblue, line width=0.9pt, -{Latex[length=2.2mm,width=1.6mm]}},
  figvecres/.style={figred, line width=1.3pt, -{Latex[length=2.6mm,width=1.9mm]}},
  figvecaux/.style={figgray, line width=0.7pt, -{Latex[length=1.8mm,width=1.3mm]}},
  figaxis/.style={figgray, line width=0.6pt, -{Latex[length=2mm,width=1.4mm]}},
  % cargas puntuales
  figcarga/.style={circle, inner sep=0pt, minimum size=5.4pt, draw=none},
  figpos/.style={figcarga, fill=figred},
  figneg/.style={figcarga, fill=figblue},
  figneutra/.style={figcarga, fill=figgray},
  % etiquetas
  figlbl/.style={font=\small},
  figlblaux/.style={font=\footnotesize, figgray},
}

% ---- pgfplots: estilo base de las gráficas ----
\pgfplotsset{
  emfig/.style={
    width=11cm, height=6.5cm,
    axis lines=left,
    xlabel style={font=\small},
    ylabel style={font=\small},
    tick label style={font=\footnotesize},
    every axis plot/.append style={line width=1pt},
    grid=major,
    grid style={figgray!20},
    legend cell align=left,
    legend style={font=\footnotesize, draw=figgray!40},
  },
}

% (compilar desde el directorio de la clase para que la ruta relativa resuelva).
% ---------------------------------------------------------------------------

% ---- Paleta (alineada con las cajas del preámbulo: keybox/notebox/warnbox) ----
\definecolor{figblue}{HTML}{1F4E79}   % azul principal (líneas, keybox)
\definecolor{figred}{HTML}{C0392B}    % rojo de énfasis (warnbox)
\definecolor{figamber}{HTML}{B8860B}  % ámbar (notebox)
\definecolor{figgray}{HTML}{555555}   % auxiliares, ejes, guías

% ---- CircuiTikz: escala y grosor de los componentes ----
% Guarda \ifdefined: la electrostática (Clases 1-14) no carga `circuitikz`, y
% sin la guarda el \input revienta con `Undefined control sequence \ctikzset`.
% Cuando el curso llegue a circuitos, el bloque se activa solo.
\ifdefined\ctikzset
  \ctikzset{
    bipoles/thickness=1,
    resistors/scale=0.9,
    capacitors/scale=0.9,
  }
\fi

% ---- TikZ: estilos reutilizables ----
% Nota: las puntas se declaran como `-{Latex[...]}` y no como `->`. La forma
% con llaves NO contiene un `>` literal, así que es inmune al `>` activo que
% introduce babel español (ver CLAUDE.md §7.3.3).
\usetikzlibrary{arrows.meta,calc,angles,quotes}

\tikzset{
  figline/.style={figblue, line width=0.9pt},
  figaux/.style={figgray, dashed, line width=0.6pt},
  % vectores
  figvec/.style={figblue, line width=0.9pt, -{Latex[length=2.2mm,width=1.6mm]}},
  figvecres/.style={figred, line width=1.3pt, -{Latex[length=2.6mm,width=1.9mm]}},
  figvecaux/.style={figgray, line width=0.7pt, -{Latex[length=1.8mm,width=1.3mm]}},
  figaxis/.style={figgray, line width=0.6pt, -{Latex[length=2mm,width=1.4mm]}},
  % cargas puntuales
  figcarga/.style={circle, inner sep=0pt, minimum size=5.4pt, draw=none},
  figpos/.style={figcarga, fill=figred},
  figneg/.style={figcarga, fill=figblue},
  figneutra/.style={figcarga, fill=figgray},
  % etiquetas
  figlbl/.style={font=\small},
  figlblaux/.style={font=\footnotesize, figgray},
}

% ---- pgfplots: estilo base de las gráficas ----
\pgfplotsset{
  emfig/.style={
    width=11cm, height=6.5cm,
    axis lines=left,
    xlabel style={font=\small},
    ylabel style={font=\small},
    tick label style={font=\footnotesize},
    every axis plot/.append style={line width=1pt},
    grid=major,
    grid style={figgray!20},
    legend cell align=left,
    legend style={font=\footnotesize, draw=figgray!40},
  },
}


% ---- Hyperlinks & TOC ----
\usepackage[hidelinks]{hyperref}
\usepackage{bookmark}

% ---- Enumerate ----
\usepackage{enumitem}

% ---- Miscellaneous ----
\usepackage{icomma}
\usepackage{siunitx}
\usepackage{textcomp}
\usepackage[bottom]{footmisc}

% ---- Title ----
\title{\textbf{Clase 8: Laplace en cilíndricas y método de las imágenes} \\[0.3em]
        \large{Campo y carga inducida en la esfera conductora}}
\author{Transcripto y expandido --- Curso Electromagnetismo (OpenFING)}
\date{}

\begin{document}

\maketitle
\thispagestyle{empty}
\newpage

\tableofcontents
\newpage


\section{Cierre del problema de la esfera conductora}

Se retoma el resultado de la Clase 7: una esfera conductora descargada de radio
$a$ en un campo uniforme $E_0\hat z$, con potencial para $r > a$

\[
\phi(r,\theta) = -E_0\,r\cos\theta + E_0\,\frac{a^3}{r^2}\cos\theta
\]

donde el primer término es lo que habría sin esfera ($-E_0 z$) y el segundo es el
efecto de la polarización.

\begin{notebox}[La constante aditiva]
El potencial está definido a menos de una constante. Con la constante elegida,
$\phi = 0$ sobre la esfera; con otra elección el valor sería otro.
\end{notebox}

\subsection{El campo eléctrico}

Con simetría azimutal el campo tiene dos componentes,
$\vec E = E_r\,\hat e_r + E_\theta\,\hat e_\theta$, que salen de
$\vec E = -\grad\phi$ en esféricas:

\[
E_r = -\pdv{\phi}{r}
\qquad\qquad
E_\theta = -\frac{1}{r}\pdv{\phi}{\theta}
\]

\begin{warnbox}[Un consejo, no de física sino de vida]
Ese signo menos es \textbf{la fuente principal de errores de cálculo} en Física
III y en Electromagnetismo. La razón es psicológica: como es lo fácil, uno lo deja
para el final, hace la cuenta larga llena de derivadas y se olvida. O peor: lo
pone en el primer término de la derivada y no en el segundo.

\textbf{La receta: poner el signo primero, en factor de todo, y recién después
derivar.} No es un error conceptual, es de distracción --- y por eso es tan común.
\end{warnbox}

Derivando el primer término respecto de $r$ sale $-E_0\cos\theta$, y del segundo,
con $r^{-2}$ en juego, sale $-2E_0 a^3 r^{-3}\cos\theta$. Aplicando el signo:

\[
\boxed{\;E_r = E_0\cos\theta\left(1 + \frac{2a^3}{r^3}\right)\;}
\]

El $1$ es lo que habría sin la esfera; el otro término es producto de la
polarización.

Para la componente angular, poniendo el $-1/r$ delante y abriendo corchete, la
derivada del coseno da $-\sin\theta$ en ambos términos:

\[
E_\theta = -\frac{1}{r}\left[E_0\,r\sin\theta - E_0\frac{a^3}{r^2}\sin\theta\right]
\qquad\Longrightarrow\qquad
\boxed{\;E_\theta = -E_0\sin\theta\left(1 - \frac{a^3}{r^3}\right)\;}
\]

(Dentro de la esfera el potencial es constante, así que $\vec E = 0$.)

\subsection{Verificación en la superficie}

Ya se sabía, de la Clase 4, que el campo justo afuera de un conductor tiene que
ser \textbf{normal a la superficie}. Eso obliga a que la componente tangencial
---que acá es $E_\theta$--- se anule en $r=a$. Y efectivamente:

\[
E_\theta\big|_{r=a} = -E_0\sin\theta\left(1 - \frac{a^3}{a^3}\right) = 0
\qquad\qquad
E_r\big|_{r=a} = 3E_0\cos\theta
\]

\begin{keybox}[La solución pasa el test]
Conviene hacer siempre esta verificación: es barata y detecta un factor 2 perdido
en cualquier parte de la cuenta.
\end{keybox}

\begin{notebox}[Un consejo de parcial]
Si en el examen la cuenta no da normal y no se encuentra el error, conviene
decírselo al corrector ---«sé que esto tiene que ser normal, me equivoqué en las
cuentas pero no en la física»---. El error de cálculo y el error conceptual no
pesan igual.
\end{notebox}

\subsection{La densidad superficial de carga}

Dentro del conductor no hay carga volumétrica, pero en la superficie sí se induce
carga. Usando el resultado general de la Clase 4, $E_n = \sigma/\varepsilon_0$,
con $E_n = E_r$:

\[
\boxed{\;\sigma(\theta) = \varepsilon_0\,E_r\big|_{r=a} = 3\varepsilon_0 E_0\cos\theta\;}
\]

\textbf{No es uniforme}: depende del ángulo, y además \textbf{cambia de signo}.

\begin{table}[H]
  \centering
  \begin{tabular}{@{}ll@{}}
    \toprule
    $\theta$ & $\sigma$ \\
    \midrule
    $0 \le \theta < \pi/2$   & positiva \\
    $\theta = \pi/2$         & nula \\
    $\pi/2 < \theta \le \pi$ & negativa \\
    \bottomrule
  \end{tabular}
\end{table}

Y eso es lo que se espera físicamente: con el campo apuntando según $+\hat z$, las
cargas negativas se acumulan del lado de los $\theta$ grandes y las positivas del
lado opuesto.

\subsection{Carga total y líneas de campo}

La esfera se supuso descargada, así que la carga total debe dar cero. Vale la pena
verificarlo. Con el elemento de superficie en esféricas,
$dS = r^2\sin\theta\,d\theta\,d\varphi$, evaluado en $r=a$:

\[
Q = \int_0^{2\pi}\!\!\int_0^{\pi} \sigma(\theta)\;a^2\sin\theta\,d\theta\,d\varphi
= 2\pi a^2\int_0^{\pi} 3\varepsilon_0 E_0\cos\theta\,\sin\theta\,d\theta
\]

(la integral en $\varphi$ da $2\pi$ porque nada depende de $\varphi$). Con el
cambio $u = \cos\theta$, $du = -\sin\theta\,d\theta$, y los límites
$\theta=0\to u=1$, $\theta=\pi\to u=-1$:

\[
Q = -2\pi a^2\,3\varepsilon_0 E_0\int_{1}^{-1} u\,du = 0
\]

porque es la integral de una función \textbf{impar} sobre un intervalo centrado en
cero. No sorprende, pero vale la pena confirmarlo: es una verificación
independiente de toda la cuenta.

\begin{figure}[H]
  \centering
  % Clase 8 — el resultado del problema de la esfera, en sus dos caras:
% (a) las lineas de campo, que llegan NORMALES a la superficie (el dibujo que el
%     docente hizo en el pizarron disculpandose por su calidad);
% (b) la densidad inducida sigma(theta), que cambia de signo y da carga total nula.
\begin{tikzpicture}[scale=1]

% ---------------- (a) lineas de campo ----------------
\begin{scope}
  % lineas lejanas: no tocan la esfera
  \foreach \yy in {-1.75,1.75}{
    \draw[figvecaux] plot[smooth, tension=0.9]
      coordinates {(-2.4,\yy) (-1.0,\yy) (0,\yy) (1.0,\yy) (2.4,\yy)};
  }
  % lineas que se deforman y llegan normales
  \draw[figvecaux] plot[smooth, tension=0.85]
    coordinates {(-2.4,1.05) (-1.35,1.05) (-0.62,0.88)};
  \draw[figvecaux] plot[smooth, tension=0.85]
    coordinates {(0.62,0.88) (1.35,1.05) (2.4,1.05)};
  \draw[figvecaux] plot[smooth, tension=0.85]
    coordinates {(-2.4,-1.05) (-1.35,-1.05) (-0.62,-0.88)};
  \draw[figvecaux] plot[smooth, tension=0.85]
    coordinates {(0.62,-0.88) (1.35,-1.05) (2.4,-1.05)};
  \draw[figvecaux] plot[smooth, tension=0.85]
    coordinates {(-2.4,0.35) (-1.5,0.35) (-1.02,0.28)};
  \draw[figvecaux] plot[smooth, tension=0.85]
    coordinates {(1.02,0.28) (1.5,0.35) (2.4,0.35)};
  \draw[figvecaux] plot[smooth, tension=0.85]
    coordinates {(-2.4,-0.35) (-1.5,-0.35) (-1.02,-0.28)};
  \draw[figvecaux] plot[smooth, tension=0.85]
    coordinates {(1.02,-0.28) (1.5,-0.35) (2.4,-0.35)};
  % esfera
  \fill[figgray!25] (0,0) circle (1.06);
  \draw[figline, line width=1.1pt] (0,0) circle (1.06);
  % normales en los puntos de llegada
  \draw[figvecres] (152:1.06) -- (152:1.62);
  \draw[figvecres] (-152:1.06) -- (-152:1.62);
  \node[figlblaux, figred, fill=white, inner sep=1.2pt] at (-2.05,0.92) {normal};
  \node[figlbl, align=center] at (0,-2.65) {(a) l\'ineas de campo};
  \node[figlblaux, align=center] at (0,-3.35)
    {paralelas lejos, deformadas cerca,\\ normales al llegar};
\end{scope}

% ---------------- (b) densidad inducida ----------------
\begin{scope}[xshift=6.6cm, yshift=-0.4cm]
  \begin{axis}[emfig, width=6.4cm, height=4.6cm, anchor=center, at={(0,0)},
    xlabel={$\theta$}, ylabel={$\sigma$},
    domain=0:pi, samples=90,
    xmin=-0.18, xmax=3.5, ymin=-1.35, ymax=1.35,
    xtick={0,1.5708,3.1416}, xticklabels={$0$,$\pi/2$,$\pi$},
    ytick=\empty,
    axis lines=middle,
    xlabel style={anchor=north west, yshift=-3pt}]
    \addplot[figblue] {cos(deg(x))};
    \addplot[figred, only marks, mark=*, mark size=1.6pt, forget plot]
      coordinates {(1.5708,0)};
    \node[figlblaux, figred]  at (axis cs:0.42,0.42)  {$\sigma>0$};
    \node[figlblaux, figblue] at (axis cs:2.42,-0.62) {$\sigma<0$};
  \end{axis}
  \node[figlbl, align=center] at (0,-2.25) {(b) $\sigma=3\varepsilon_0E_0\cos\theta$};
  \node[figlblaux, align=center] at (0,-2.95)
    {cambia de signo en $\theta=\pi/2$;\\ la integral sobre la esfera da $Q=0$};
\end{scope}

\end{tikzpicture}

  \caption*{\small\itshape Las dos caras del mismo resultado: la geometría de las
  líneas y la distribución de carga que la produce. Que $\sigma$ sea impar
  alrededor de $\pi/2$ es lo que hace que la carga total se anule.}
\end{figure}

\textbf{Las líneas de campo.} Lejos de la esfera son paralelas y equidistribuidas,
porque el campo es uniforme. Al acercarse se deforman, y \textbf{llegan normales a
la superficie}. Las más alejadas del eje pasan sin tocar la esfera.

\subsection{Sobre la superposición}

Una pregunta de clase que vale la pena registrar: ¿por qué no se resolvió el
problema por superposición desde el principio?

\textbf{Sí vale superposición} ---el campo resultante es la suma del campo uniforme
más el campo producido por las cargas inducidas, que resulta ser
\textbf{estrictamente dipolar}---. Lo que pasa es que \textbf{a priori no se sabía
cuál era esa distribución de carga}: $\sigma$ es \emph{inducida por el propio
campo externo}, y se obtuvo recién al final de la cuenta.

\begin{keybox}
La superposición no ayudaba porque le faltaba un dato. Lo que el método sí
garantizó es que la solución hallada, al cumplir Laplace y todas las condiciones
de borde, es \emph{la} solución, por unicidad.
\end{keybox}

\section{Laplace en coordenadas cilíndricas}

Se repite todo el programa en otras coordenadas, para mostrar que el método de
separación de variables es general.

\subsection{Planteo y separación}

Igual que en esféricas no se trató el caso más general ---se pidió simetría
azimutal---, acá se pide \textbf{invariancia por traslaciones en $z$}: el problema
es \textbf{plano}.

\[
\phi = \phi(r,\theta), \qquad \pdv{\phi}{z} = 0
\]

\begin{notebox}[El ejemplo]
Un \textbf{cilindro conductor descargado muy largo}, de radio $a$, con su eje
perpendicular al plano del dibujo, colocado en un campo uniforme. El dibujo es
idéntico al de la esfera, \textbf{pero el objeto es completamente distinto}. La
hipótesis de «muy largo» es la que elimina la dependencia en $z$.
\end{notebox}

Con el último término del laplaciano cilíndrico anulado, y proponiendo
$\phi(r,\theta) = Y(r)\,S(\theta)$:

\[
\frac{S(\theta)}{r}\dv{r}\big(r\,Y'(r)\big) + \frac{Y(r)}{r^2}\,S''(\theta) = 0
\]

Multiplicando por $r^2/(Y S)$ se separan las variables:

\[
\frac{r}{Y(r)}\dv{r}\big(r\,Y'(r)\big) = -\frac{S''(\theta)}{S(\theta)} = K
\]

y otra vez cada miembro \textbf{no depende de ninguna variable} ---el izquierdo no
depende de $\theta$ y está igualado a algo que no depende de $r$---, así que ambos
son la misma \textbf{constante de separación} $K$.

\subsection{La ecuación angular y la periodicidad}

\[
S''(\theta) + K\,S(\theta) = 0
\]

Ésta sí es «facilonga»: es la ecuación del resorte. Pero \textbf{cuál de sus
familias de soluciones sirve lo decide una condición física}.

\begin{keybox}[La restricción]
Se consideran problemas en que $\theta$ recorre \textbf{todo} el intervalo
$[0,2\pi]$. Como al dar la vuelta entera se vuelve al mismo punto del espacio,
$S(\theta)$ tiene que ser una \textbf{función periódica de período $2\pi$}.
\end{keybox}

\begin{warnbox}[Un contraejemplo que aparece en el práctico]
Dos planos semi-infinitos que forman una cuña de ángulo $\alpha$, con $\phi=0$ en
uno y $\phi=\phi_0$ en el otro, y se busca el potencial entre ellos. El problema se
escribe naturalmente en cilíndricas y adentro vale Laplace, \textbf{pero $\theta$
sólo recorre $[0,\alpha]$}: no se puede dar la vuelta, no hay periodicidad, y
\textbf{el análisis que sigue no se le aplica}. Ese caso tiene condiciones de borde
distintas y \textbf{más soluciones}, y hay que rehacer el análisis.
\end{warnbox}

Con esa restricción:

\begin{itemize}[leftmargin=*]
  \item \textbf{Si $K < 0$}, las soluciones son \textbf{exponenciales reales}, y
        ninguna combinación de exponenciales reales es periódica.
        \textbf{Se descartan.} Luego $K \ge 0$.
  \item \textbf{Si $K = \omega^2 \ge 0$}, la ecuación es $S'' + \omega^2 S = 0$ y
        $S(\theta) = A\cos(\omega\theta) + B\sin(\omega\theta)$.
\end{itemize}

Para que el período sea exactamente $2\pi$, $\omega$ tiene que ser
\textbf{entero}.

\[
\boxed{\;K = n^2, \qquad n \in \mathbb{Z}\;}
\]

\begin{notebox}[La misma estructura, mucho menos trabajo]
Comparado con el camino esférico ---serie de potencias, recurrencia, polinomios de
Legendre, criterio de convergencia--- este fue mucho más corto. Pero el argumento
es el mismo: la condición física sobre el dominio es la que \textbf{cuantiza} la
constante de separación.
\end{notebox}

\subsection{La ecuación radial}

Con $K = n^2$: \quad $r\,\dv{r}\big(r\,Y'(r)\big) = n^2\,Y(r)$. \quad Otra vez
lineal, de segundo orden, con coeficientes dependientes de $r$.

\begin{notebox}[Ejercicio con la misma pista de la clase anterior]
Cambio de variable $v = \ln r$.
\end{notebox}

Hay que separar dos casos:

\begin{itemize}[leftmargin=*]
  \item \textbf{$n \ge 1$:} \quad $Y(r) = A\,r^{-n} + B\,r^{\,n}$
  \item \textbf{$n = 0$:} la ecuación se reduce a $\dv{r}(r\,Y') = 0$, de donde
        $r\,Y' = C$, $Y' = C/r$ y \quad $Y(r) = C\ln r + D$
\end{itemize}

\begin{notebox}[Asimetría entre los dos casos]
El caso $n=0$ da \textbf{una sola} solución angular (la constante, con
$\omega=0$) en lugar de dos, pero \textbf{dos} radiales: el logaritmo y la
constante. Reaparece el potencial logarítmico de la línea infinita de carga que ya
había salido en la Clase 6.
\end{notebox}

\subsection{La base de soluciones}

\[
\boxed{\;\phi(r,\theta) = C\ln r + D
+ \sum_{n=1}^{\infty}\Big(a_n r^{\,n} + b_n r^{-n}\Big)
  \Big(A_n\cos n\theta + B_n\sin n\theta\Big)\;}
\]

\begin{notebox}[Hay una redundancia deliberada en la escritura]
Por cada $n$ aparecen cuatro constantes, pero sólo \textbf{tres} son
independientes, porque un factor común de la primera pareja se puede absorber en
la segunda. Se deja así porque «el que puede más puede menos»: al imponer
condiciones de borde algunas constantes sólo aparecerán como productos.
\end{notebox}

\textbf{Ejercicio dejado (y que está en el práctico):} resolver el cilindro
conductor descargado en el campo uniforme con este desarrollo, imponiendo las
condiciones de borde igual que con la esfera.

\begin{warnbox}[Una diferencia delicada con el caso esférico]
Para cerrar el problema de la esfera se usó el desarrollo multipolar y el hecho de
que la carga total era nula. Acá \textbf{eso no se puede aplicar directamente}: el
desarrollo multipolar vale para distribuciones \textbf{localizadas en una región
finita}, y un cilindro infinito no lo es. Se puede estar muy lejos comparado con
el radio, pero no comparado con la longitud. Al final el problema se resuelve
igual, pero el argumento requiere más cuidado.
\end{warnbox}

\section{El método de las imágenes}

La motivación vuelve al \textbf{teorema de unicidad} de la Clase 5: si se
encuentra \emph{una} solución que cumple todas las condiciones de borde, ganó ---
«todos los golpes están permitidos». La consecuencia práctica es que conviene
coleccionar métodos: cuantos más golpes se sepan, más duro se le puede pegar a la
ecuación.

\subsection{Carga frente a un plano conductor}

\textbf{El problema:} un plano conductor infinito y una carga $q$ a distancia $d$
de él. Se quiere resolver Laplace \textbf{sólo del lado de la carga}; lo que pasa
del otro lado no interesa.

\begin{keybox}[La idea rectora: un plano conductor es un espejo]
De hecho muchos espejos se construyen así. Un observador parado del lado de la
carga, si no se da cuenta de que eso es un espejo, \textbf{ve dos cargas}.
Entonces la solución de este lado tiene que ser la misma que la del problema de dos
cargas.
\end{keybox}

\begin{figure}[H]
  \centering
  % Clase 8 — el metodo de las imagenes, caso plano. La carga imagen es FICTICIA y
% la solucion vale SOLO del lado de la carga real: del otro lado del espejo es
% otra historia.
% Las dos cotas `d` van bien abajo y separadas: a la altura del eje se pisan entre
% si y con los rotulos de las cargas.
\begin{tikzpicture}[scale=1]

  % region donde NO vale la solucion
  \fill[figgray!12] (-3.3,-2.6) rectangle (0,2.2);
  \draw[figline, line width=1.4pt] (0,-2.6) -- (0,2.2);
  \node[figlbl, figblue] at (0,2.5) {plano conductor};
  \node[figlblaux, align=center] at (-2.55,-1.35) {aqu\'i la soluci\'on\\ \emph{no} vale};

  \draw[figaxis] (-3.3,0) -- (3.5,0);
  \node[figlblaux] at (3.7,0) {$x$};

  % carga real
  \node[figpos] at (1.7,0) {};
  \node[figlbl, figred] at (1.7,0.42) {$+q$};

  % carga imagen
  \node[figneg] at (-1.7,0) {};
  \node[figlbl, figblue] at (-1.7,0.42) {$q'=-q$};
  \node[figlblaux, align=center] at (-1.7,1.28) {imagen\\ (ficticia)};

  % carga inducida real sobre el plano
  \foreach \yy in {-1.2,-0.85,0.85,1.2}{\node[figneg] at (0,\yy) {};}
  \node[figlblaux, figblue, fill=white, inner sep=1.2pt] at (0.85,1.55) {$\sigma<0$};

  % cotas, bien abajo
  \draw[figaux] (-1.7,-0.22) -- (-1.7,-2.15);
  \draw[figaux] (1.7,-0.22) -- (1.7,-2.15);
  \draw[figaux] (0,-1.45) -- (0,-2.15);
  \draw[figvecaux] (0,-1.95) -- (1.7,-1.95);
  \draw[figvecaux] (1.7,-1.95) -- (0,-1.95);
  \node[figlbl, fill=white, inner sep=1.5pt] at (0.85,-1.95) {$d$};
  \draw[figvecaux] (0,-1.95) -- (-1.7,-1.95);
  \draw[figvecaux] (-1.7,-1.95) -- (0,-1.95);
  \node[figlbl, fill=figgray!12, inner sep=1.5pt] at (-0.85,-1.95) {$d$};

  \node[figlbl, align=center] at (0.1,-3.35)
    {En $x=0$ los dos denominadores coinciden:
     $\phi\propto\dfrac{q+q'}{\sqrt{d^{2}+y^{2}+z^{2}}}$,};
  \node[figlbl, align=center] at (0.1,-4.25)
    {constante s\'olo si $q'=-q$, y entonces $\phi_0=0$.};

\end{tikzpicture}

  \caption*{\small\itshape La carga imagen no existe: es un artificio que
  reproduce la condición de borde. Lo que sí existe es la densidad negativa
  inducida sobre el plano.}
\end{figure}

Se propone una \textbf{carga imagen} $q'$, ficticia, del otro lado del plano y a
la misma distancia. Con el plano en $x=0$ y la carga real en $x=d$:

\[
\phi(x,y,z) = \frac{1}{4\pi\varepsilon_0}\left[
\frac{q}{\sqrt{(x-d)^2+y^2+z^2}} + \frac{q'}{\sqrt{(x+d)^2+y^2+z^2}}\right]
\]

\textbf{Imponer la condición de borde.} Sobre el plano, $x=0$, los dos
denominadores se vuelven \textbf{idénticos}:

\[
\phi(0,y,z) = \frac{1}{4\pi\varepsilon_0}\cdot\frac{q + q'}{\sqrt{d^2+y^2+z^2}}
\]

Para que eso sea \textbf{constante} ---independiente de $y$ y $z$--- hay una sola
posibilidad:

\[
\boxed{\;q' = -q \qquad\text{y entonces}\qquad \phi_0 = 0\;}
\]

\[
\phi(x,y,z) = \frac{q}{4\pi\varepsilon_0}\left[
\frac{1}{\sqrt{(x-d)^2+y^2+z^2}} - \frac{1}{\sqrt{(x+d)^2+y^2+z^2}}\right]
\]

\begin{notebox}[Se verifican todas las condiciones]
Es constante sobre el plano; cerca de la carga real domina el término de la carga
y se comporta como debe; y $\phi\to0$ a distancia grande, coherente con que la
carga está localizada.
\end{notebox}

\begin{warnbox}[Ésta es la solución sólo de este lado del espejo]
Del otro lado es otra historia --- si el conductor es macizo, allí el potencial es
constante e igual a cero.
\end{warnbox}

La carga inducida es \textbf{negativa}, y es razonable: una carga positiva atrae
las negativas hacia la zona cercana; en un conductor infinito las positivas quedan
a lo lejos.

\begin{notebox}[Comparación con la alternativa]
El mismo problema se podría atacar con simetría azimutal y polinomios de Legendre,
pero sería «una pesadilla»: hay una serie infinita, un punto singular donde no vale
Laplace, y habría que resolver por regiones y empalmar. Acá fueron dos cargas.
\end{notebox}

\subsection{Densidad inducida sobre el plano}

Sobre el plano el campo tiene que ser normal, es decir según $\hat x$. Derivando
(y poniendo el signo menos \textbf{primero}, como manda §1.1):

\[
E_x = -\pdv{\phi}{x}
= \frac{q}{4\pi\varepsilon_0}\left[
\frac{x-d}{\big((x-d)^2+y^2+z^2\big)^{3/2}}
- \frac{x+d}{\big((x+d)^2+y^2+z^2\big)^{3/2}}\right]
\]

Evaluando en $x=0$, los dos denominadores coinciden y los numeradores se suman:

\[
\boxed{\;\sigma(y,z) = \varepsilon_0 E_x\big|_{x=0}
= -\,\frac{q\,d}{2\pi\big(d^2+y^2+z^2\big)^{3/2}}\;}
\]

Nótese que $\varepsilon_0$ se cancela, como debe.

\begin{notebox}[Ejercicio dejado en clase]
Verificar que la carga inducida total es exactamente la opuesta a la carga real,
\[
\int_{-\infty}^{\infty}\!\!\int_{-\infty}^{\infty}\sigma\,dy\,dz = -q
\]
\textbf{Truco:} pasar a coordenadas polares en el plano. El resultado tiene que dar
$-q$: la carga distribuida sobre el plano es equivalente a la carga imagen.
\end{notebox}

\subsection{Dos espejos perpendiculares}

\textbf{El problema:} dos planos conductores perpendiculares y una carga $q$ en el
cuadrante que forman. Se resuelve Laplace en ese cuadrante. \textbf{Hacen falta
tres imágenes}, no una ni dos.

\begin{figure}[H]
  \centering
  % Clase 8 — dos espejos perpendiculares: hacen falta TRES imagenes, no una ni dos.
% La cuarta carga (la imagen de las imagenes) es la que restaura la constancia en
% el primer plano despues de agregar la imagen del segundo.
\begin{tikzpicture}[scale=1]

  \fill[figblue!8] (0,0) rectangle (2.6,2.6);

  \draw[figline, line width=1.4pt] (-2.6,0) -- (2.6,0);
  \draw[figline, line width=1.4pt] (0,-2.6) -- (0,2.6);
  \node[figlblaux] at (2.9,-0.22) {$x$};
  \node[figlblaux] at (-0.22,2.85) {$y$};

  % las cuatro cargas
  \node[figpos] at (1.3,1.55) {};
  \node[figlbl, figred] at (1.75,1.78) {$+q$};
  \node[figneg] at (1.3,-1.55) {};
  \node[figlbl, figblue] at (1.78,-1.78) {$-q$};
  \node[figneg] at (-1.3,1.55) {};
  \node[figlbl, figblue] at (-1.78,1.78) {$-q$};
  \node[figpos] at (-1.3,-1.55) {};
  \node[figlbl, figred] at (-1.78,-1.78) {$+q$};

  % las reflexiones, punteadas
  \draw[figaux] (1.3,1.33) -- (1.3,-1.33);
  \draw[figaux] (1.08,1.55) -- (-1.08,1.55);
  \draw[figaux] (-1.3,1.33) -- (-1.3,-1.33);
  \draw[figaux] (-1.08,-1.55) -- (1.08,-1.55);

  \node[figlblaux, align=center, fill=figblue!8, inner sep=1.5pt] at (1.95,0.62)
    {regi\'on donde\\ se resuelve};

  \node[figlbl, align=center] at (0,-3.35)
    {Con dos espejos en \'angulo recto se ven \textbf{cuatro} im\'agenes,
     y eso es exactamente\\[0.2em]
     lo que hace falta: agregar s\'olo una rompe la constancia en el otro plano.};

\end{tikzpicture}

  \caption*{\small\itshape La construcción matemática reproduce exactamente lo que
  hace la óptica: con dos espejos en ángulo recto se ven cuatro imágenes.}
\end{figure}

\begin{table}[H]
  \centering
  \begin{tabular}{@{}ll@{}}
    \toprule
    \textbf{Posición} & \textbf{Carga} \\
    \midrule
    reflejada en el plano horizontal & $-q$ \\
    reflejada en el plano vertical   & $-q$ \\
    reflejada en ambos (la diagonal) & $+q$ \\
    \bottomrule
  \end{tabular}
\end{table}

\textbf{Por qué la cuarta es indispensable.} La carga real más su imagen en un
plano hacen que \emph{ese} plano sea equipotencial. Pero al agregar la imagen del
otro plano, esa condición se rompe. Para restaurarla hay que agregar \textbf{la
imagen de la imagen}. Y entonces, por simetría, las cuatro cargas juntas dejan
equipotenciales \textbf{los dos} planos a la vez.

\subsection{Carga frente a una esfera conductora}

El ejemplo serio. Una \textbf{esfera conductora} de radio $a$ centrada en el origen
y una carga $+q$ sobre el eje, a distancia $d > a$ del centro. Se busca el
potencial \textbf{afuera} de la esfera, con la condición $\phi\big|_{r=a} = 0$.

\begin{notebox}[La intuición del espejo curvo]
Un espejo curvo también da una imagen, pero \textbf{deformada}. Se espera entonces
que la imagen de una carga puntual siga siendo una carga puntual, pero con
\textbf{otra posición} $d'$ y \textbf{otro valor} $q'$ --- ninguno de los dos obvio
de antemano.

Lo que sí se sabe por simetría es que la imagen está \textbf{sobre el eje}: el
problema es invariante por rotación alrededor de él, así que nada puede
desimetrizarla.
\end{notebox}

Con la carga imagen en $x = d'$:

\[
\phi(x,y,z) = \frac{1}{4\pi\varepsilon_0}\left[
\frac{q}{\sqrt{(x-d)^2+y^2+z^2}} + \frac{q'}{\sqrt{(x-d')^2+y^2+z^2}}\right]
\]

\textbf{Imponer $\phi=0$ sobre la esfera.} Desarrollando los cuadrados aparece la
combinación $x^2+y^2+z^2$, que sobre la esfera vale $a^2$ ---«parece que estuviera
hecho a propósito»---:

\[
\frac{q}{\sqrt{a^2+d^2-2dx}} = -\,\frac{q'}{\sqrt{a^2+d'^2-2d'x}}
\qquad\text{para todo } x
\]

Elevando al cuadrado y multiplicando cruzado:

\[
q^2\big(a^2+d'^2-2d'x\big) = q'^2\big(a^2+d^2-2dx\big)
\]

Como esto debe valer \textbf{para todo $x$} sobre la esfera, se igualan por
separado la parte constante y la parte lineal en $x$ --- un sistema de dos
ecuaciones con dos incógnitas:

\[
q^2\big(a^2+d'^2\big) = q'^2\big(a^2+d^2\big)
\qquad\qquad
q^2\,d' = q'^2\,d
\]

De la segunda, $q'^2 = q^2\,d'/d$. Sustituyendo en la primera, $q^2$ se
simplifica:

\[
a^2+d'^2 = \frac{d'}{d}\big(a^2+d^2\big)
\qquad\Longrightarrow\qquad
d'^2 - \frac{a^2+d^2}{d}\,d' + a^2 = 0
\]

Resolviendo la cuadrática y \textbf{completando el cuadrado} bajo la raíz:

\[
d' = \frac{\big(a^2+d^2\big) \pm \sqrt{\big(a^2+d^2\big)^2 - 4a^2d^2}}{2d}
= \frac{\big(a^2+d^2\big) \pm \big|a^2-d^2\big|}{2d}
\]

Como $d > a$, el valor absoluto vale $d^2-a^2$, y quedan dos candidatos:

\begin{table}[H]
  \centering
  \begin{tabular}{@{}ccl@{}}
    \toprule
    \textbf{Signo} & \textbf{Valor} & \textbf{¿Sirve?} \\
    \midrule
    $+$ & $\dfrac{2d^2}{2d} = d$ & \textbf{No}: es la posición de la carga original \\[0.7em]
    $-$ & $\dfrac{2a^2}{2d} = \dfrac{a^2}{d}$ & \textbf{Sí} \\
    \bottomrule
  \end{tabular}
\end{table}

Sustituyendo en $q'^2 = q^2 d'/d = q^2 a^2/d^2$ y eligiendo el signo negativo (la
imagen tiene que compensar a la carga real, que está del lado cercano a esa parte
de la superficie):

\[
\boxed{\;d' = \frac{a^2}{d}\;}
\qquad\qquad
\boxed{\;q' = -\,q\,\frac{a}{d}\;}
\]

\begin{figure}[H]
  \centering
  % Clase 8 — el espejo curvo: carga q frente a una esfera conductora. La imagen
% sigue siendo puntual pero cambia de posicion Y de valor, y para el caso
% DESCARGADO hace falta una segunda imagen, en el centro.
\begin{tikzpicture}[scale=1]

  % esfera conductora
  \fill[figgray!25] (0,0) circle (1.15);
  \draw[figline, line width=1.2pt] (0,0) circle (1.15);

  \draw[figaxis] (-2.0,0) -- (4.4,0);
  \node[figlblaux] at (4.6,0) {$x$};

  % radio a
  \draw[figvecaux] (0,0) -- (118:1.15);
  \node[figlblaux, fill=figgray!25, inner sep=1.2pt] at (-0.32,0.68) {$a$};

  % carga real
  \node[figpos] at (3.3,0) {};
  \node[figlbl, figred] at (3.3,0.44) {$+q$};

  % carga imagen, dentro de la esfera
  \node[figneg] at (0.68,0) {};
  \node[figlbl, figblue, fill=figgray!25, inner sep=1.2pt] at (0.78,0.46) {$q'$};

  % la segunda imagen, en el centro
  \node[figpos] at (0,0) {};
  \draw[figvecaux] (-1.42,1.55) -- (-0.16,0.14);
  \node[figlbl, align=center, fill=white, inner sep=2pt] at (-1.9,2.05)
    {2.\textsuperscript{a} imagen: $+q\,a/d$\\ \emph{en el centro}};

  % cotas, abajo y separadas
  \draw[figaux] (0,-0.22) -- (0,-1.85);
  \draw[figaux] (0.68,-0.22) -- (0.68,-1.15);
  \draw[figaux] (3.3,-0.22) -- (3.3,-1.85);
  \draw[figvecaux] (0,-0.95) -- (0.68,-0.95);
  \node[figlblaux, fill=white, inner sep=1.2pt] at (0.34,-0.72) {$d'$};
  \draw[figvecaux] (0,-1.65) -- (3.3,-1.65);
  \draw[figvecaux] (3.3,-1.65) -- (0,-1.65);
  \node[figlbl, fill=white, inner sep=1.5pt] at (1.65,-1.65) {$d$};

  % resultados
  \node[figlbl] at (1.2,-2.65)
    {$\boxed{\,d'=\dfrac{a^{2}}{d}\,}\qquad
      \boxed{\,q'=-\,q\dfrac{a}{d}\,}$};
  \node[figlblaux, align=center] at (1.2,-3.85)
    {Como $d>a$ resulta $d'<a$: la imagen cae \emph{dentro} de la esfera,\\[0.2em]
     fuera de la regi\'on donde se resuelve el problema.};
  \node[figlblaux, align=center] at (1.2,-5.15)
    {Con s\'olo $q'$ la esfera queda \emph{cargada}; la carga del centro la
     descarga\\[0.2em] sin romper la equipotencial (la lleva a otra constante,
     ya no cero).};

\end{tikzpicture}

  \caption*{\small\itshape La imagen cae dentro de la esfera ---fuera de la región
  donde se resuelve---, y hace falta una segunda para pasar del caso cargado al
  descargado.}
\end{figure}

\begin{notebox}[Y es consistente]
Como $d>a$, resulta $d' < a$: la carga imagen queda \textbf{dentro} de la esfera,
que es donde tiene que estar --- fuera de la región donde se resuelve el problema.
\end{notebox}

\textbf{Pero ésa no es la solución que interesa.} La configuración hallada
corresponde a una esfera con \textbf{carga neta} $q' = -qa/d$: el potencial cero en
la superficie exige que la esfera esté cargada.

\textbf{El caso realista es la esfera descargada.} Para pasar a él hay que agregar
una \textbf{segunda carga imagen} que (1) no rompa la condición de equipotencial y
(2) cancele la carga neta.

\begin{keybox}[¿Dónde ponerla?]
No en un lugar simétrico ---eso rompería la primera condición---. La respuesta es
\textbf{en el centro}: una carga puntual en el centro de la esfera mantiene la
superficie equipotencial (a otro valor constante, ya no cero), que es todo lo que
se pide.
\[
q'' = +\,q\,\frac{a}{d} \quad\text{en el origen}
\qquad\Longrightarrow\qquad
q' + q'' = 0
\]
\end{keybox}

\begin{notebox}[La moraleja del método]
En palabras del docente: no es un procedimiento científico sino
\textbf{artístico}. No hay una regla que dé la solución; hay que ser astuto para
elegir dónde poner las cargas, y sólo funciona en geometrías suficientemente
simples como para saber cuál es la imagen. Pero con experiencia se le agarra la
mano --- y el teorema de unicidad garantiza que, si las condiciones de borde se
cumplen, la solución hallada es \emph{la} solución.
\end{notebox}

\bigskip
\noindent\emph{Continúa en la Clase 9 ---cuya transcripción no está publicada en
OpenFING--- con la ecuación de Laplace en coordenadas cilíndricas y el método de
las imágenes aplicados a más casos, y el paso a los medios dieléctricos.}

\end{document}
